\section{软件连接与分层架构}
\label{chapter:02}

% block-id: auto.02.p1
\hypertarget{auto.02.p1}{}
系统分为研究工作流、执行适配和物理模型三个层次。工作流定义任务输入、操作顺序、检查依据和完成条件；适配层管理 Wolfram、AEDT 及数值后端的连接和执行；模型层保存公式、离散实现与路线测试。参数与产物通过共同任务接口流转，使文本教程、命令行和 MCP 入口能够引用同一份已保存输入。

% block-id: auto.02.p2
\hypertarget{auto.02.p2}{}
Wolfram 接入固定 AgentTools 2.2.0，来源提交为 9c99e2da5286c2ff7ac2298956107cd9334c6e90；既有验证记录使用 Wolfram 内核 15.0.0。工具提供代码执行、测试以及 Notebook 读写。数值表达式返回和 TestReport 支持当前计算检查。归档的 rc.2 前端案例分别检查了 Notebook 的打开、保存、重开与图形回读。\cite{AGENTTOOLS}

% block-id: auto.02.p3
\hypertarget{auto.02.p3}{}
Ansys 接入固定 PyAEDT-MCP 0.2.dev0，来源提交为 80ae7c3cb47e5de7b744a2898061510299d3bcef；建模适配采用 PyAEDT 1.8.0。服务器在独立 CPython 环境中运行，项目运行器启动并记录所属 AEDT Student gRPC 会话，再由 MCP 连接到该会话。AEDT 内置脚本环境与外部 CPython 具有不同的模块和对象接口，调用记录因此保留实际执行位置、工具名与返回值。\cite{PYAEDTMCP,PYAEDT,ANSYSCOMMONMCP}

% block-id: auto.02.p4
\hypertarget{auto.02.p4}{}
默认配置采用本地标准输入输出传输，配置生成器也提供本机 HTTP 方式。安装材料固定上游来源、依赖和许可，连接诊断核对实际路径与工具清单。Wolfram、AEDT 及字体是运行环境的组成部分，版本和许可证条件随运行记录保存。其他 MCP 客户端可使用相同任务说明和协议接口，本文端到端记录对应 Codex 与标准 MCP 客户端桥接。

% block-id: auto.02.p5
\hypertarget{auto.02.p5}{}
研究源码与软件连接之间保留清晰边界，使模型能够通过命令行独立运行，也便于定位失败发生在哪一层。客户端会话中产生的临时代码在正式验收前写入文件；运行清单记录文件哈希和参数。由此，后续修订能够通过代码差异和新的运行身份追溯到原始尝试。
